\documentclass[10pt,a4paper]{amsart}
\usepackage[margin=19mm]{geometry}
\usepackage{amsmath,amssymb,mathtools}
\usepackage[scr=rsfs,cal=euler]{mathalfa}
\usepackage{xfrac}
\usepackage[unicode,hidelinks,bookmarks=false]{hyperref}
\newcommand{\ie}{i.e.\ }
\newcommand{\cf}{cf.\ }
\newcommand{\resp}{respectively\ }
\newcommand{\Subsection}{\subsection}
\providecommand{\nobreakdash}{\nobreak\hbox{-}}
\setlength{\emergencystretch}{2em}
\multlinegap0pt
\usepackage{amssymb}
\usepackage{tikz}
\usepackage{enumerate}
\usepackage[all]{xy}
\newtheorem{prop}{Proposition}
\newcommand{\Aa}{\mathcal{A}}
\newcommand{\Mm}{\mathcal{M}}
\newcommand{\RR}{\mathbb{R}}
\newcommand{\dR}{\mathrm{dR}}
\title{On the real homotopy type of compact complex surfaces}
\author{Joana Cirici}
\date{Source: arXiv:2601.01958v1 (CC BY 4.0)}
\begin{document}
\maketitle

\noindent\emph{Source and licence.} On the real homotopy type of compact complex surfaces. Version arXiv:2601.01958v1 (CC BY 4.0), distributed by arXiv under the Creative Commons Attribution 4.0 International licence, http://creativecommons.org/licenses/by/4.0/, as recorded in the arXiv licence metadata for this exact version and shown on the abstract page https://arxiv.org/abs/2601.01958v1. Full licence notice: \texttt{licenses/arxiv-2601.01958-CC-BY-4.0.txt}.\par\smallskip

\noindent\textit{Editorial scope.} Counted unit: the article's prop prop (source0.tex lines 800--808). The article's complete original proof of it is delivered with it (source0.tex lines 809--813, one \texttt{proof} environment of 5 lines). No mathematics is written, repaired or re-derived: the statement and the proof are the article's own lines, in the article's own notation, and no retained line is substituted or re-flowed.  Retained uncounted as the source-local context the proof needs: lines 713--715.  Nothing was omitted from any retained span, so the package carries no omission record.  The retained lines cite 2 works, whose entries are transcribed verbatim from the article's own bibliography source in the e-print and delivered with short editorial labels recorded in the item's reference mapping.  No retained line contains a figure environment, an image-inclusion command or drawing code, so no image is delivered and the package declares no asset dependency.  Editorial formatting only, in addition: an AMS \texttt{amsart} preamble that restates the article's own declarations and macros used by the retained lines, namely the environments \texttt{prop} on separate counters, together with the standard packages those lines need.  The retained environments print in this document as Proposition 1 in order of appearance, whereas the article prints the same items with the numbering of its own full document.  The complete native source of the article as distributed in the arXiv e-print for this exact version is bundled unchanged as \texttt{sources/192124/source0.tex}.
\begin{prop}\label{formalpure}
 If a (commutative) dg-algebra admits a pure weight decomposition, then it is formal. In particular, all higher Massey products vanish.
\end{prop}

\begin{prop}Assume that $\Aa_{\dR}(X)$ admits a 1-model $\Mm$ together with a weight decomposition.
\begin{enumerate}
 \item If the weight decomposition is pure, then the real Malcev completion of $\pi_1(X)$ is determined
 by the cup product $H^1(X)\times H^1(X)\to H^2(X)$ and its Malcev Lie algebra is quadratically presented.
 \item If the induced weights in cohomology satisfy
 $H^1(X)\cong H^1_1\oplus H^1_2$ and $H^2(X)\cong H^2_2\oplus H^2_3$
 then the real Malcev completion of $\pi_1(X)$ is determined by the bigraded nilpotent Lie algebra of $(\pi_1(X)/\Gamma_4)\otimes\RR$.
\end{enumerate}
\end{prop}

\begin{proof}
First, note that weight decompositions carry on to minimal models so we may assume that such 1-model is minimal (see for instance Lemma 3.2 of \cite{Thom}).
In the pure case, the minimal 1-model is formal by Proposition \ref{formalpure} and hence may be computed from the cup product $H^1(X)\times H^1(X)\to H^2(X)$. The second statement follows from the fact that
the associated graded Lie algebra has generators with weights -1 and -2 (corresponding to the negative values of the weights in $H^1$) and relations of weights -2 and -3 (corresponding to the negative values of the weights in $H^2$). Consequently, the graded Lie algebra modulo its fourth order commutators reconstructs the full tower. We refer to \cite{Morgan} for more details.
\end{proof}

\begin{thebibliography}{99}
\bibitem[Morgan]{Morgan} J.W. Morgan. \newblock The algebraic topology of smooth algebraic varieties. \newblock {\em Inst. Hautes \'Etudes Sci. Publ. Math.}, (48):137--204, 1978.
\bibitem[Thom]{Thom} U.~Buijs, F.~Cantero~Mor\'an, and J.~Cirici. \newblock Weight decompositions of {T}hom spaces of vector bundles in rational homotopy. \newblock {\em J. Homotopy Relat. Struct.}, 15(1):1--26, 2020.
\end{thebibliography}
\end{document}
